\section{Population-Foster closure of the 432-pair active-invariant
orbit
gap}

\textbf{Proof-first standalone theorem, 2026-08-12 PDT.} This note
closes the support-pair orbit which had been removed by an invariant
positive on only two coordinates before species permutations were taken.
The original active-coordinate invariant is not used as a globally
proper Foster function. Instead, the exact orbit splits into
deficiency-zero pairs, pairs with a statewise factorial fourth-power
drift, and two explicit rank-two exceptional families with a direct
population-level service proof.

Finite computation is used only for support-orbit, deficiency,
invariant, tier, affine-feasibility, and orbit-representative
identities. No orientation, rate vector, population box, or stochastic
history is enumerated.

\subsection{1. The family, process, and one population
potential}

Begin with the inherited ordered shielded/available seed pairs. Among
them, 110 were assigned to the common-active-invariant prebranch. Close
these pairs under all species permutations and linkage reversal, then
subtract the orbit of every other inherited seed pair. Denote the
resulting exclusive family by \(\mathcal G\).

Fix \(P\in\mathcal G\), an arbitrary strongly connected directed graph
on each of its two linkage supports, arbitrary positive labelled rate
constants, and a closed irreducible population class \(\Gamma\). Let
\(X_t\) be the ordinary stochastic mass-action chain. Choose \(K>1\) and
put

\[
 G(x)=K+\sum_{i=A,B,C}\log(x_i!),
 \qquad W(x)=G(x)^4.                                  \tag{1.1}
\]

This is one population function: there is no reaction mark,
chart-dependent linear correction, or entrance recharge. It is
nonnegative and proper on \(\mathbb N_0^3\).

The chain is nonexplosive for every pair considered here. A quadratic
source cannot increase total molecularity because every target has
molecularity at most two. Every population-increasing reaction therefore
has source degree zero or one, total increasing intensity
\(O(1+|x|_1)\), and bounded jump size. A linear pure-birth comparison
bounds total population on finite time intervals; each population
sublevel is finite and has bounded total hazard.

\subsection{2. Exact finite
decomposition}

The authoritative finite certificate verifies all of the following.

\begin{enumerate}
\def\labelenumi{\arabic{enumi}.}
\tightlist
\item
  The 110 seeds have a 714-pair orbit. Removing the orbit of the other
  5,059 inherited seeds leaves exactly \(432\) pairs in \(\mathcal G\).
\item
  Every pair has stoichiometric rank two. Its primitive normal has
  exactly one zero coefficient and two coefficients of the same strict
  sign. Each choice of zero coordinate occurs 144 times.
\item
  Exactly 174 pairs have full-network deficiency zero.
\item
  Of the remaining 258 pairs, 234 have no affine-feasible failure of the
  corrected S-tier-superlevel cut.
\item
  The last 24 pairs form two disjoint
  species-permutation/linkage-reversal orbits, each of size 12,
  represented by \[
  \begin{array}{c|c|c}
    &L_1&L_2\\ \hline
  \mathrm I&\{A,A+B\}&\{2A,2C,A+C\}\\
  \mathrm {II}&\{A,A+B\}&\{C,2A,B+C\}.
  \end{array}                                        \tag{2.1}
  \] Their primitive nonnegative invariants are respectively \(A+C\) and
  \(A+2C\). Thus the only unbounded coordinate on a fixed class is
  \(B\), the zero coordinate of the invariant.
\item
  The 192 affine-feasible corrected-cut failures include 48
  deficiency-zero pairs and the 24 pairs in (2.1). In every failed row,
  the invariant-zero coordinate is exactly the unique active coordinate.
\end{enumerate}

The principal exact fingerprints are

\begin{CodeBlock}
all 432 pairs
5516d6071b2b9d07b0e4e02613b9caee217ba3ebb0082e21f2bc664e6247ea36
deficiency-zero 174
d894c296de7c24e3c855f116228b852ef940b15ac913556ffea9c87e6655a974
non-DZ / no feasible failure 234
4c6420ad47c2ee57736484b213595dde69c3734ce20d19ff4f97aeb31ab77d6a
exceptional 24
051a641f3987ec93b129ad044d96292a97e536e9ef3d2724234dc4af9bfdef69
type-I orbit 12
50804d58a48bb1a2683014442a436761ec0ba0df34af3ef3ac0d1939fe850886
type-II orbit 12
cc76d3c7bcc7942f956f1efc8cf0f718ad1e6e08d911f8f8127ebc3cf8c5002f
invariant manifest
9dec8108276e9d439c18aacda1ec35d9bac08e097f8833e3446c50b40d8148ca
failure/invariant alignment
a9368dd934b7ac6135c3df4866e2322700d3a607dd58f8327aeb709065880ab2
\end{CodeBlock}

These are support identities only. The analytic implications are proved
in Sections 4--6.

The authoritative certificate bytes are

\begin{CodeBlock}
src/active_invariant_orbit_gap_432_certificate.py
31fa24a20e18546e9c623d3aaf6d3b845c1708d5782f86333c02417fa366cd53

tests/test_active_invariant_orbit_gap_432_certificate.py
09c434fc162ff33f51e331e298ec2e35407a8e0501f1a3b8d771bf33c2fe708b
\end{CodeBlock}

Its nine dedicated tests and the combined 22-test replay pass.

\subsection{3. Elementary fourth-power
estimates}

On any fixed conserved \((A,C)\)-phase, as \(B\to\infty\),

\[
             G(A,B,C)=B\log B+O(B).                  \tag{3.1}
\]

If a bounded reaction changes \(G\) by \(u\), then exactly

\[
 (G+u)^4-G^4=4G^3u+6G^2u^2+4Gu^3+u^4.              \tag{3.2}
\]

Consequently, uniformly over a finite \((A,C)\)-phase, for all
sufficiently large \(B\):

\[
\begin{array}{rcll}
 u&=&-\log B+O(1)&\Longrightarrow\quad
          \Delta W\le-cG^3\log B,\\
 u&=& \log(B+1)+O(1)&\Longrightarrow\quad
          (\Delta W)^+\le CG^3\log B,\\
 |u|&\le&C&\Longrightarrow\quad |\Delta W|\le CG^3.
\end{array}                                          \tag{3.3}
\]

All constants in (3.1)--(3.3) may depend on the fixed population class
and rate vector. No uniformity across different classes or rates is
needed.

We will also use the following elementary random-time criterion.

\begin{quote}
\textbf{Lemma 3.1 (bounded-depth population Foster rule).} Let a
nonexplosive CTMC have a proper nonnegative function \(V\) and a finite
set \(K_0\). Suppose that from every state outside \(K_0\) there is an
all-clock stopping time \(\tau>0\), ending after between one and
\(d<\infty\) actual jumps, such that \[
      \mathbb E_x[V(X_\tau)-V(x)+\tau]\le-1,        \tag{3.4}
\] and \(V(X_\tau)+\tau\) is integrable. Then the mean hitting time of
\(K_0\) is finite. On a closed irreducible class the chain is positive
recurrent.
\end{quote}

\subsubsection{Proof}

Let \(\sigma_0=0\), and recursively complete the selected episode from
its actual endpoint. Define \(N\) to be the first positive episode count
\(j\) for which the path segment from \(\sigma_{j-1}\) through
\(\sigma_j\) visits \(K_0\); complete that episode only for accounting.
Conditional summation of (3.4) through \(N\wedge n\) gives

\[
 \mathbb E V(X_{\sigma_{N\wedge n}})
 +\mathbb E\sigma_{N\wedge n}
 +\mathbb E(N\wedge n)\le V(x).                     \tag{3.5}
\]

Fatou and monotone convergence give \(\mathbb EN<\infty\) and finite
mean accounting time. Indeed, on \(\{N=\infty\}\), the episode endpoints
include an infinite subsequence of actual jump times because every
episode contains at least one jump. Nonexplosion forces
\(\sigma_n\to\infty\), contradicting the finite expectation obtained
from (3.5). Thus \(N<\infty\) almost surely, and the actual hit of
\(K_0\) occurs no later than \(\sigma_N\).

For the recurrence promotion, fix \(o\in K_0\cap\Gamma\), which is
nonempty by the hitting result. From each of the finitely many
\(c\in K_0\cap\Gamma\), irreducibility gives a finite actual labelled
path to \(o\). On the event that each prescribed label fires before
every competitor, that path is completed. Across the finitely many path
states, this event has a positive uniform probability and the attempt
has finite uniform mean duration. The possible first-competitor
endpoints form a finite set; the preceding hitting bound gives a finite
uniform mean return from them to \(K_0\cap\Gamma\). Geometric repetition
gives a finite mean hit of \(o\). A singleton class is already positive
recurrent. Otherwise, after the first jump from \(o\), the same bound
gives a finite mean return to \(o\). Hence the closed irreducible class
is positive recurrent. \(\square\)

\subsection{4. The 174 deficiency-zero
pairs}

Each directed linkage graph is strongly connected, so the full network
is weakly reversible. For every positive rate vector, weak reversibility
and full deficiency zero give a positive complex-balanced vector
\(c\in(0,\infty)^3\). On a closed irreducible class \(\Gamma\), the
stochastic product form is

\[
 \pi_\Gamma(x)=Z_\Gamma^{-1}
       \prod_{i=A,B,C}{c_i^{x_i}\over x_i!},
 \qquad x\in\Gamma.                                  \tag{4.1}
\]

It is normalizable because its mass is bounded by \(\exp(c_A+c_B+c_C)\),
the full-lattice product-Poisson mass before the constant factors.
Nonexplosion was proved in Section 1. Thus every closed irreducible
class is positive recurrent. This is the class-restricted
Anderson--Craciun--Kurtz deficiency-zero theorem, applied to the full
two-linkage network, not separately to its linkages.

\subsection{5. The 234 pairs with no feasible corrected-cut
failure}

We prove a direct generator theorem under the population potential
(1.1). Suppose, to the contrary, that no finite set supports
\(\mathcal LW\le-1\) on a fixed class \(\Gamma\). Properness supplies a
divergent sequence \(x_n\in\Gamma\) with \(\mathcal LW(x_n)>-1\).
Extract one of the finitely many exact tier/cap descriptors.

Because the sequence lies in one affine stoichiometric class, the
necessity direction of the exact affine-feasibility theorem makes this
descriptor feasible. By the defining 234-pair identity it therefore
satisfies the corrected S-tier-superlevel cut. For every strong
orientation, that cut gives a D-descending reaction sourced in the
literal global top S-tier.

Let \(A_n\) be the largest enabled source propensity and let
\(g_n\to\infty\) be its certified logarithmic D-gap. The exact factorial
calculation gives

\[
 \mathcal LG(x_n)\le-cA_ng_n,
 \qquad \sum_r a_r(x_n)\le CA_n,
 \qquad |\Delta_rG(x_n)|\le C\log(2+|x_n|_1).        \tag{5.1}
\]

Using (3.2) reaction by reaction,

\[
\begin{split}
 \mathcal LW={}&4G^3\mathcal LG
 +6G^2\sum_ra_r(\Delta_rG)^2\\
 &+4G\sum_ra_r(\Delta_rG)^3
 +\sum_ra_r(\Delta_rG)^4.
\end{split}                                          \tag{5.2}
\]

Since \(G(x_n)\asymp |x_n|_1\log(2+|x_n|_1)\), every remainder in (5.2)
is lower order than \(-G^3A_ng_n\). Hence
\(\mathcal LW(x_n)\to-\infty\), a contradiction. Therefore

\[
                         \mathcal LW\le-1            \tag{5.3}
\]

outside a finite subset of \(\Gamma\). Continuous-time Foster, localized
at finite \(W\)-sublevels and followed by Fatou, gives finite mean hit
of that finite set and positive recurrence. This branch has no chart
exits and no random-potential seam.

\subsection{6. The two exceptional
orbits}

Relabel an exceptional pair as in (2.1). The invariant makes \(A,C\)
range over a finite set on every fixed class, while \(B\) is the only
coordinate that can diverge.

\subsubsection{6.1 Type II: pointwise
service}

Fix the invariant value

\[
                         M=A+2C.                     \tag{6.1}
\]

If \(M=0\), all sources are disabled and each population state is an
absorbing singleton. Assume \(M>0\). Since both linkage graphs are
strong, the source \(A+B\) has a positive total rate to \(A\), and the
source \(B+C\) has a positive total rate to the B-free complexes \(C\)
or \(2A\). Thus the aggregate rate of B-removing reactions is

\[
       B(\delta_1A+\delta_2C)\ge c_M B.              \tag{6.2}
\]

Every such reaction changes \(G\) by \(-\log B+O_M(1)\). All B-creating
sources are among \(A,C,2A\); their aggregate propensity is bounded by a
classwise constant. Reactions which do not change B also have bounded
propensity and bounded \(G\)-increment. Equations (3.3) and (6.2) give

\[
 \mathcal LW
 \le-c_M B G^3\log B+C_MG^3\log B+C_MG^3
 \longrightarrow-\infty.                            \tag{6.3}
\]

Hence \(\mathcal LW\le-1\) outside a finite set, and ordinary
continuous-time Foster proves positive recurrence.

\subsubsection{6.2 Type I: activation followed by
service}

Now fix

\[
                          M=A+C.                     \tag{6.4}
\]

The first linkage necessarily contains both arrows
\(A\rightleftarrows A+B\). Whenever \(A>0\), the B-death rate is
\(\delta AB\), whereas its B-birth rate is \(\beta A=O_M(1)\). The
second linkage has no B-coordinate and has total rate \(O_M(1)\).
Therefore, at a state \(z=(A,B,C)\) with \(A>0\),

\[
 \mathcal LW(z)
       \le-c_M B G(z)^3\log B+C_MG(z)^3\log B.      \tag{6.5}
\]

Its total hazard satisfies

\[
                   c_M B\le\lambda(z)\le C_M(1+B).  \tag{6.6}
\]

If \(J\) is the next actual jump and \(T_J\) its holding time, then
(6.5)--(6.6) imply, for large \(B\),

\[
 \mathbb E_z[W(X_J)-W(z)+T_J]
       ={\mathcal LW(z)+1\over\lambda(z)}
       \le-c'_MG(z)^3\log B.                         \tag{6.7}
\]

It remains to cover \(A=0\). If \(M=0\), or if \(M=1\) and \(C=1\), no
complex in either linkage is enabled; the state is an absorbing
singleton. If \(M\ge2\), the only enabled source is \(2C\). Every
outgoing target from \(2C\) is either \(2A\) or \(A+C\), because its
linkage is strong. Thus the first actual jump occurs at a fixed positive
classwise rate, leaves B unchanged, and lands at a state \(z\) with
\(A>0\). Its \(G\)-increment is bounded by a classwise constant, and
hence its \(W\)-increment is at most \(C_MG^3\). Append the next actual
jump from \(z\). By (6.7),

\[
 \mathbb E_{(0,B,M)}
 [W(X_\tau)-W(0,B,M)+\tau]
 \le C_MG^3-c'_MG^3\log B+C_M\le-1                 \tag{6.8}
\]

for all sufficiently large \(B\). From \(A>0\), define \(\tau\) to be
one actual jump; from \(A=0,M\ge2\), define it to be the two jumps in
(6.8). Every clock is retained. The endpoint changes B by at most one,
\(A,C\) stay in their finite invariant phase, and the duration has
finite mean bounded uniformly in B. Thus the endpoint and duration in
(6.8) are integrable.

The complement of the large-B region is finite on the fixed class. Lemma
3.1 applied to (6.7)--(6.8) proves positive recurrence of every
non-singleton type-I class.

\subsection{7. Pair theorem and seam
disposition}

\begin{quote}
\textbf{Theorem 7.1 (the exclusive 432-pair orbit).} For every support
pair in \(\mathcal G\), every strongly connected orientation on its two
linkage supports, every fixed positive labelled rate vector, and every
closed irreducible population class, the stochastic mass-action chain is
nonexplosive and positive recurrent.
\end{quote}

\subsubsection{Proof}

The exact decomposition in Section 2 is disjoint and exhaustive. Section
4 handles 174 pairs, Section 5 handles 234, and Section 6 handles the
final 24. The counts sum to 432. \(\square\)

This proof intentionally supersedes the earlier marked- \(W\)
drift-or-chart-exit draft. A terminal strongly connected component can
cycle entirely through chart exits, so common marked potential plus
unweighted exit circulation is not sufficient. Here the 234 passing
pairs have statewise population-generator drift, type II has statewise
direct service, and type I has an unconditional one- or two-jump
population Foster episode from every large state. No chart exit, cap
promotion, marked recharge, or weighted terminal seam remains.

\subsection{8. Frozen analytic dependencies and claim
boundary}

The passing calculation uses the exact analytic statements in:

\begin{CodeBlock}
research_notes/s_tier_superlevel_cut_and_affine151_corrected.md
d91f369d34cadfb28ddb872df8fb9f6d17799ec207da29933037f55ae95f0407

research_notes/stoichiometric_gate_feasibility.md
27b40b61903ae6c2e223d007ec08323ec9aec10e9198deb99d2d7c60d878d007

research_notes/universal_fourth_power_one_active_interface.md
9d4239f4fc6b45a9522b94b09523c9f98ac7a3b089c919bd9594f12409c78cc2
\end{CodeBlock}

The deficiency-zero publication scope is independently checked in the
\href{proof_first_publication_primary_literature_audit.md}{primary-literature
audit}.

The theorem is pair-level. A final global union must pin the
authoritative 432 certificate and prove that this orbit and the other
certified branches cover the intended global support universe with the
stated precedence.
